\documentclass[11pt,a4paper]{article}
\usepackage[margin=25mm]{geometry}
\usepackage{amsmath,amssymb,amsthm,hyperref}
\newtheorem{lemma}{Lemma}
\newtheorem{proposition}[lemma]{Proposition}
\newcommand{\gen}[1]{\langle#1\rangle}
\title{Supported intervals and the one-point parity extension}
\author{Human Compression project: bounded L044 proof component}
\date{}
\begin{document}\maketitle
Let $G=\gen{\rho_0,\ldots,\rho_{r-1}}\leq A_n$ be a
transitive string C-group, with all maximal omissions intransitive.
Suppose the unique eligible $i$-edge joins $a\in A$ to $b\in B$,
the two $G_i$-orbits, and $0<i<r-1$. Thus every generator
other than $\rho_i,\rho_{i-1},\rho_{i+1}$ fixes both $a,b$.
Put $L=\gen{\rho_j:j<i}$ and $R=\gen{\rho_j:j>i}$.
They commute and $G_i=LR$.

\begin{lemma}[The actual two supported intervals]
Suppose, among the distinguished indices other than $i$, those
acting nontrivially on $A$ form one interval,
and those acting nontrivially on $B$ also form one interval.
Then, after interchanging $A,B$, $L$ is supported on and
transitive on $A$, and $R$ is supported on and transitive on $B$.
Their markings are faithful even string C-markings, their
maximal omissions are intransitive, and $a,b$ are endpoint anchors.
If $|A|,|B|\geq3$, their noncommuting diagrams are connected.
\end{lemma}
\begin{proof}
Each active interval lies on one side of $i$. It must include
the adjacent index on that side: otherwise every generator
would fix the orbit's crossing endpoint, contradicting
transitivity on that orbit. They cannot both lie on the same
side, because then every generator on the other, nonempty
side would fix both $A$ and $B$ pointwise and would be the
identity. Hence the intervals occupy opposite sides. Every
index on its side belongs to its interval, since a missing
index would again act trivially on both orbits.
This proves the exact support and faithful-action assertions.
Evenness, the string property and IP follow by restriction
of these unchanged disjointly supported submarkings.

If a maximal omission in $L$ were transitive on $A$, adjoining
$R$, transitive on $B$, and $\rho_i$, joining the two orbits,
would make the corresponding global omission transitive.
That is forbidden. The same argument works for $R$.
Fixation of $a,b$ by all nonendpoint members is the
unique-crossing property.

To see connectedness, suppose two nontrivial commuting marked
components of $L$ existed. Every component except the one
containing its last index fixes $a$. Commutation implies that
it fixes the whole orbit of $a$ under the other components,
which is $A$, since $L$ is transitive. Its faithful action would
then be trivial. Thus no such component exists. The proof for
$R$ is identical.
\end{proof}

\begin{lemma}[A faithful even one-point extension]
In the supported situation of Lemma~1, suppose
$L=A_A$ and $R=A_B$, with $|A|,|B|\geq5$.
At least one of the two interval groups has a one-point
extension which is a faithful even string C-group of rank
one larger, transitive on the enlarged support, with every
maximal omission intransitive and an endpoint anchor.
\end{lemma}
\begin{proof}
Write $s=\rho_i$. Its only cross-orbit transposition is
$(a\ b)$, so uniquely
\[
                 s=(a\ b)u v,
\]
where $u$ is an involution supported on $A\setminus\{a\}$
and $v$ is one supported on $B\setminus\{b\}$, allowing
identity factors. As $s$ is even,
$\operatorname{sgn}(u)\operatorname{sgn}(v)=-1$.
Interchange the sides if necessary so that $u$ is odd and
$v$ is even. Then $v\in R=A_B$.

Let $P=\gen{L,s}$ and project its action to the invariant set
$X=A\cup\{b\}$. The image $\Phi$ is generated by the unchanged
marking of $L$ together with $t=(a\ b)u$, an even involution.
Every element of $P$ acts on $B\setminus\{b\}$ either as the
identity or as $v$. Therefore the projection kernel is either
trivial or $\gen v$. In the latter nontrivial case $v\in P\cap R$,
contrary to the full intersection property of $G$: the index
sets defining $P$ and $R$ are disjoint. Hence projection is
injective. This proves the complete marked isomorphism, not
merely an equality of abstract orders, and transfers full IP.

The image is transitive on $X$ because $L$ is transitive on
$A$ and $t$ joins $a,b$. Its new last endpoint has anchor $b$.
If an omission of an old member were transitive on $X$, the
same actual subgroup of the corresponding global omission,
together with $R$ on $B$, would make that global omission
transitive. Omitting $t$ leaves $b$ fixed. Thus every omission
of $\Phi$ is intransitive. Its diagram is connected: $L$ is
connected by Lemma~1, and $t$ cannot centralize all of $L$,
since it moves a point of its transitive support outside that
support. String commutation then connects it to the adjacent
endpoint. All stated properties follow.
\end{proof}

\begin{proposition}[The large two-interval branch of the rank induction]
Assume the supported-interval hypotheses above, with
$m_A=|A|\geq12$ and $m_B=|B|\geq12$.
Assume the alternating-group rank bound
$d\leq\lfloor(m-1)/2\rfloor$ for every smaller degree
$12\leq m<n$. Then
\[
                       r\leq\left\lfloor\frac{n-1}{2}\right\rfloor.
\]
\end{proposition}
\begin{proof}
Every interval action has the anchor and omission hypotheses
of Lemma~1. If it is a proper primitive subgroup of its
alternating group, the accepted complete primitive bound gives
rank at most $\lfloor(m-3)/2\rfloor$. If it is imprimitive,
the accepted anchored imprimitive theorem gives rank at most
$\lfloor(m-2)/2\rfloor$. If it is the full alternating group,
the stated induction hypothesis gives rank at most
$\lfloor(m-1)/2\rfloor$. No unanchored form of source
Proposition~5.13 is used.

Let $d_A,d_B$ be the interval ranks, so $r=d_A+d_B+1$.
If at least one interval group is proper in its alternating
group, the bounds just proved give
\[
 r\leq\frac{m_A+m_B-3}{2}+1=\frac{n-1}{2}.
\]
Otherwise both are full alternating groups. Apply Lemma~2,
choosing the side $A$ after relabelling. The enlarged group
$\Phi\leq A_{m_A+1}$ has rank $d_A+1$, a connected diagram,
an anchor, and all omissions intransitive. Its degree satisfies
$12\leq m_A+1<n$. The same primitive/imprimitive/alternating
trichotomy therefore gives
\[
          d_A+1\leq\left\lfloor\frac{m_A}{2}\right\rfloor,
  \qquad d_B\leq\left\lfloor\frac{m_B-1}{2}\right\rfloor.
\]
Adding proves the assertion, with all induction calls at
strictly smaller degrees.
\end{proof}

\paragraph{Scope and attribution.}
The source is Cameron--Fernandes--Leemans--Mixer,
\emph{Highest rank of a polytope for $A_n$},
arXiv:1605.09173v1. This is the large-degree part of its Section~5.2,
Propositions~5.15--5.17. The one-point projection is the
source's sesqui-extension idea, written with an explicit
kernel and full marked-IP proof. Its essential use of the
full ambient IP is retained. The endpoint anchor in the
smaller action is proved, including in the new projected
extension; it is not inferred merely from having a quotient.
The accepted complete proper-primitive and anchored-imprimitive
proofs are paid dependencies. The alternating step is an
explicit strong-induction hypothesis, not an already completed
proof of L044. Small supported orbits, their Table~2 models,
the remaining mixed-support cases and the equality-graph
branch remain outside this component. No new computational,
formal, original or question credit is claimed.
\end{document}
